% LaTeX-ready fragment for the shared project report.
% Required packages: amsmath, booktabs, tabularx, graphicx.
% Bibliography entries are provided in report_references_cnn.bib.
% If the main file is not compiled from the repository root, redefine
% \reportfigdir to the copied figure directory before inputting this fragment.
\providecommand{\reportfigdir}{docs/figures}

\section{Dataset}
\label{sec:dataset}

The project uses VisDrone2019-DET, an aerial object detection benchmark
collected from drone-mounted cameras under different cities, viewpoints,
altitudes, weather conditions, and scene densities \cite{zhu2018vision,du2019visdrone}.
The task contains ten foreground categories: pedestrian, people, bicycle, car,
van, truck, tricycle, awning-tricycle, bus, and motor. In contrast to
street-level datasets, VisDrone contains many small, densely packed, and
partially occluded objects. These properties directly influenced the
preprocessing, detector configuration, and evaluation protocol.

\subsection{Exploratory Data Analysis (EDA)}
\label{subsec:dataset-insight}

The exploratory analysis was executed reproducibly with random seed 42 in the
project EDA notebook. Dataset and box statistics include only annotations with
\texttt{status=keep}, exactly matching objects exported for training; rejected
rows are audited separately. Image-quality results use a fixed sample of 250
images per split and are interpreted as diagnostics rather than full-population
estimates.

After ignored and invalid annotations were removed, the data used in this
project contained 8,629 images and 457,065 valid objects. Table
\ref{tab:split-statistics} summarizes the three processed splits.

\begin{table}[htbp]
  \centering
  \caption{VisDrone split statistics after preprocessing.}
  \label{tab:split-statistics}
  \resizebox{\textwidth}{!}{%
  \begin{tabular}{lrrrrrr}
    \toprule
    Split & Images & Raw rows & Retained boxes & Mean/image & Median/image & Maximum/image \\
    \midrule
    Train      & 6,471 & 353,550 & 343,204 & 53.04 & 42 & 902 \\
    Validation &   548 &  40,169 &  38,759 & 70.73 & 65 & 317 \\
    Test-dev   & 1,610 &  77,547 &  75,102 & 46.65 & 36 & 461 \\
    \midrule
    Total      & 8,629 & 471,266 & 457,065 & -- & -- & -- \\
    \bottomrule
  \end{tabular}}
\end{table}

The integrity audit found no missing image--annotation pair, unreadable image,
or malformed annotation row. Of 471,266 raw rows, 457,065 (96.99\%) are valid
training objects, 14,200 are ignored regions with \texttt{score=0}, and one
training box has non-positive size. No image becomes empty after filtering.
Thus, the labeled data pass the basic integrity checks, while ignored regions
must remain separate from ordinary negative background. No retained box
required boundary clipping under the current validation rules.

\paragraph{Class Distribution.}
Table \ref{tab:class-distribution} reports retained instances across all three
labeled splits. Cars and pedestrians constitute 64.80\% of all objects, while
awning-tricycles constitute 0.96\%. In training, the maximum-to-minimum
instance ratio is 44.6:1. The instance-level loss therefore receives many more
updates from common classes. Although mAP weights categories equally, it still
hides which rare classes fail; per-class AP and a confusion matrix are needed.

\begin{table}[htbp]
  \centering
  \caption{Class distribution across the labeled VisDrone splits.}
  \label{tab:class-distribution}
  \begin{tabular}{lrr@{\hspace{1.2cm}}lrr}
    \toprule
    Class & Instances & Share & Class & Instances & Share \\
    \midrule
    Car        & 187,004 & 40.91\% & Truck            & 16,284 & 3.56\% \\
    Pedestrian & 109,187 & 23.89\% & Bicycle          & 13,069 & 2.86\% \\
    Motor      &  40,378 &  8.83\% & Bus              &  9,117 & 1.99\% \\
    People     &  38,560 &  8.44\% & Tricycle         &  6,387 & 1.40\% \\
    Van        &  32,702 &  7.15\% & Awning-tricycle  &  4,377 & 0.96\% \\
    \bottomrule
  \end{tabular}
\end{table}

\paragraph{Object Size Distribution.}
Following the COCO area convention, 60.47\% of training boxes and 68.56\% of
validation boxes are small ($\text{area}<32^2$ pixels). Table
\ref{tab:box-scale} shows that the median training object is only 25 by 27
pixels and occupies approximately 0.046\% of the image. A localization error
of only a few pixels can therefore substantially decrease IoU.

\begin{table}[htbp]
  \centering
  \caption{Bounding-box scale statistics using the COCO area convention.}
  \label{tab:box-scale}
  \begin{tabular}{lrrrrrr}
    \toprule
    Split & Med. width & Med. height & Med. area & Small & Medium & Large \\
    \midrule
    Train      & 25 px & 27 px & 680 px$^2$ & 60.47\% & 34.00\% & 5.53\% \\
    Validation & 20 px & 25 px & 520 px$^2$ & 68.56\% & 28.68\% & 2.76\% \\
    Test-dev   & 20 px & 23 px & 480 px$^2$ & 67.66\% & 29.13\% & 3.21\% \\
    \bottomrule
  \end{tabular}
\end{table}

Across all 457,065 retained boxes, the median geometry is 24$\times$26 pixels
with median area 627 pixels$^2$. A total of 38.91\% have at least one side
below 16 pixels, 72.13\% have a side below 32 pixels, and 69.19\% occupy less
than 0.1\% of their image. This relative-area diagnostic is distinct from the
COCO categories. Tiny objects affect every class: 92.44\% of people, 88.88\%
of pedestrians, 83.96\% of motors, and 81.38\% of bicycles occupy less than
0.1\% of the image. Even bus and truck contain 30.89\% and 38.12\% tiny boxes.

\begin{figure}[htbp]
  \centering
  \includegraphics[width=\textwidth]{\reportfigdir/dataset_statistics.pdf}
  \caption{Class imbalance and COCO-scale distribution. Cars dominate the
  retained labels, while small boxes account for more than 60\% of every split.}
  \label{fig:dataset-statistics}
\end{figure}

\begin{figure}[htbp]
  \centering
  \includegraphics[width=\textwidth]{\reportfigdir/class_difficulty.pdf}
  \caption{Per-class scale and visibility difficulty. ``Tiny'' denotes a box
  whose area is below 0.1\% of its image and is an EDA diagnostic rather than
  a COCO evaluation category.}
  \label{fig:class-difficulty}
\end{figure}

\paragraph{Image properties and split shift.}
Images have variable resolutions, including 1,400$\times$1,050,
1,400$\times$788, 2,000$\times$1,500, 1,360$\times$765,
1,916$\times$1,078, and 1,920$\times$1,080. Training width ranges from 480 to
2,000 pixels and its mean resolution is approximately 1,520$\times$1,002.
Validation is smaller on average (1,291$\times$726) and almost entirely has
aspect ratio 1.78, whereas training mixes ratios near 1.33 and 1.78. Resizing
to fixed independent width and height would therefore introduce geometric
distortion and potentially a split-dependent artifact.

Class composition also changes. Relative to training, validation contains 5.34
percentage points more people and 3.97 more motors, but 5.92 fewer cars.
Class-distribution Jensen--Shannon divergence is 0.0164 for train--validation
and 0.0084 for train--test-dev. The larger shift lies in density:
Kolmogorov--Smirnov distance is 0.253 for train--validation and 0.347 for
validation--test-dev.

\begin{figure}[htbp]
  \centering
  \includegraphics[width=\textwidth]{\reportfigdir/eda_diagnostics.pdf}
  \caption{Key diagnostics from the EDA notebook: image density, tiny-object
  indicators, class-composition shift, and simulated object scale after
  letterboxing.}
  \label{fig:eda-diagnostics}
\end{figure}

\paragraph{Objects per Image and Object Density.}
Validation contains 70.73 retained objects per image on average, compared with
53.04 in training and 46.65 in test-dev. Furthermore, 19.53\% of validation
images and 10.86\% of training images contain more than 100 objects; 22 images
across the labeled data contain more than 300. Proposal and output limits that
are suitable for a general-purpose dataset can therefore cap recall.

Across all valid boxes, approximately 49.5\% are unoccluded, 40.9\% are
partially occluded, 9.6\% are heavily occluded, and 4.0\% are truncated.
Motor has the highest heavy-occlusion rate (15.92\%); bicycle, tricycle, and
awning-tricycle are near 13\%. More than 71\% of people and bicycle instances
are occluded to some degree. Bus has the highest truncation rate (9.63\%),
followed by truck (7.61\%). These labels should be retained for training and
used for stratified error analysis.

\paragraph{Bounding-box Characteristics.}
Across all 457,065 retained boxes, median width is 24 pixels, median height is
26 pixels, and median area is 627 pixels$^2$. The 10th--90th percentile ranges
are 7--81 pixels for width, 10--72 pixels for height, and 90--5,200 pixels$^2$
for area. Aspect ratio, defined as width divided by height, has median 0.89 and
spans 0.41--2.08 between the 10th and 90th percentiles. The task therefore
contains portrait-like pedestrians, near-square vehicles, and wide objects in
addition to substantial scale variation.

\begin{figure}[htbp]
  \centering
  \includegraphics[width=\textwidth]{\reportfigdir/density_bbox_geometry.pdf}
  \caption{Object density and bounding-box geometry. The panels show complete
  objects-per-image distributions, log-scale area distributions with COCO
  thresholds, and width/height aspect ratios.}
  \label{fig:density-bbox-geometry}
\end{figure}

\paragraph{Qualitative Analysis.}
A fixed six-image audit gallery complements the aggregate statistics. It
contains tiny distant targets, exceptionally crowded traffic, severe
occlusion, low and high illumination, and a visibly different object
scale/viewpoint. Red rectangles denote retained ground-truth boxes. The same
gallery should later be used to compare predictions from every model family,
including representative failures rather than only successful examples.
VisDrone does not provide per-image altitude metadata, so the final panel shows
apparent scale/viewpoint variation rather than verified flight altitude.

\begin{figure}[htbp]
  \centering
  \includegraphics[width=\textwidth]{\reportfigdir/qualitative_eda_gallery.pdf}
  \caption{Six representative VisDrone conditions for qualitative audit:
  tiny objects, crowded traffic, heavy occlusion, illumination variation, and
  scale/viewpoint variation. This is a dataset audit, not a prediction result.}
  \label{fig:qualitative-eda-gallery}
\end{figure}

\paragraph{Spatial context and co-occurrence.}
Median normalized horizontal box centers lie between 0.488 and 0.504 for all
classes, indicating no strong left--right bias. Vertical position is
class-dependent: large vehicles are higher in the image (median normalized
$y\approx0.34$--0.35), while people and two-wheel vehicles are lower
($y\approx0.40$--0.46). Bus and truck have the highest boundary-contact rates,
7.01\% and 5.99\%. This drone-perspective structure supports mild bbox-safe
translation but argues against extreme crops or rotations.

The strongest image-level co-occurrences are car--van (Jaccard 0.790),
pedestrian--car (0.771), and people--motor (0.718). Context can help recognition
but may create shortcuts. Oversampling a complete image for a rare category can
also oversample common co-occurring classes, so targeted instance copy-paste is
more controllable when placement remains plausible.

\paragraph{Duplicate, leakage, and image-quality audit.}
SHA-256 found three byte-identical groups within individual splits: two in
training and one in validation. Their annotations differ and require manual
inspection. No exact duplicate crosses a split. However, 24 filename sequence
prefixes occur in both training and validation, and three same-sequence pairs
have perceptual dHash distance at most six. This does not prove direct leakage,
but it indicates temporal or scene correlation. The official split remains
necessary for benchmark comparison; a sequence-aware evaluation is preferable
when measuring generalization to unseen scenes.

On a fixed sample of 250 images per split, validation is brighter on average
(0.4518 on a normalized 0--1 scale) than training (0.3693) and test-dev
(0.3447). Mean contrast is similar and validation sharpness is slightly lower.
Because this is a sampled diagnostic, it supports mild brightness, contrast,
color, and blur augmentation but not aggressive global correction.

\paragraph{Task-level implications from EDA.}
The notebook translated these findings into controlled simulations:
The square-letterbox simulation estimates scale sensitivity across model
families. It is a dataset diagnostic rather than the exact runtime preprocessing
of any one detector.
\begin{itemize}
  \item With square letterboxing, 46.63\%, 27.63\%, and 16.65\% of boxes have
  short side below 8 pixels at inputs 640, 960, and 1,280 respectively.
  Resolution must therefore be tested as an accuracy--memory ablation.
  \item At simulated input 960, large default anchor shapes cover substantially
  fewer ground-truth shapes than an 8--128-pixel set at shape-IoU $\geq0.5$.
  This indicates that all model families must preserve fine-scale features or
  adapt their assignment mechanism to tiny targets; the implementation remains
  architecture-specific.
  \item A 640-pixel tile with 20\% overlap geometrically contains 99.81\% of
  boxes in at least one tile but costs approximately 6.59 tile inferences per
  image. Tiling is therefore a general accuracy--latency hypothesis that must
  be validated separately for each detector.
  \item Repeat-Factor Sampling with threshold 0.1 assigns every class a factor
  of 1.0 because even awning-tricycle occurs in 17.79\% of training images.
  Complete-image oversampling therefore does not solve instance imbalance.
\end{itemize}

These simulations provide hypotheses, while trained ablations determine the
final configuration. Overall, the EDA identifies tiny objects, crowded scenes,
class imbalance, occlusion, split shift, and sequence correlation as the main
data risks. It supports aspect-ratio-preserving resizing, high-resolution
experiments, increased proposal/detection limits, per-class and scale-aware
metrics, and a fixed qualitative audit set containing crowded, tiny, blurred,
and rare-class examples.

\subsection{Preprocess}
\label{subsec:preprocess}

Preprocessing was performed once before the processed dataset was uploaded to
Kaggle. The same image splits, class mapping, and COCO ground truth are shared
by the CNN-, YOLO-, and Transformer-based experiments, preventing data
conversion differences from confounding their comparison. Each raw
annotation follows

\begin{center}
\texttt{x, y, width, height, score, category, truncation, occlusion}.
\end{center}

The pipeline matched images and annotations by filename, verified image
readability, removed rows with \texttt{score <= 0}, and retained only category
IDs 1--10. Category 0 denotes ignored regions, while category 11 denotes
``others''; neither was used as a positive training label. Non-positive and
fully out-of-image boxes were discarded. Boundary-crossing boxes were clipped
to the image, after which boxes smaller than 1$\times$1 pixel were removed.
Valid labels were exported as COCO JSON with pixel-space \texttt{xywh} boxes,
area, category ID, and \texttt{iscrowd=0}. YOLO labels were exported in parallel
for the other model family in the project.

Preprocessing removed 10,345 ignored training rows and one non-positive box,
1,410 ignored validation rows, and 2,445 ignored test-dev rows. No missing
image, missing annotation, unreadable image, or malformed annotation was found.
Images were hard-linked into the processed tree, with copying as a fallback.
No offline resizing, normalization, or augmentation was performed. Original
pixels are therefore shared across architectures. Input sizing, normalization,
tensor layout, and label conversion are applied at runtime by each model family
and are documented in its experiment setup. Ignored regions are excluded from the processed
training and validation COCO files but retained in the raw annotations. The
unified comparison uses the processed ground truth only. A separate
official-compatible VisDrone evaluation may use the raw annotations to suppress
detections whose area overlaps an ignored region by at least 50\%, as specified
by the VisDrone toolkit \cite{visdrone_toolkit}; those values are not mixed with
the unified results.

\subsection{Dataset Augmentation}
\label{subsec:augmentation}

The EDA defines common augmentation requirements for the complete detection
task, independent of downstream architecture. Augmentation is applied online
and only to training samples, so original images remain unchanged and no
augmented copies are stored. Validation and test images are not stochastically
transformed, ensuring that all model families are compared on a fixed
distribution.

Any geometric operation must transform the image and its bounding boxes with
the same mapping, clip boxes to the image boundary, and remove invalid remnants
without desynchronizing their class labels. Since many VisDrone targets are
only a few pixels wide, aggressive cropping, blur, rotation, or downscaling can
erase the target rather than create useful invariance. The EDA therefore
supports horizontal flipping, mild affine and photometric perturbations, and
carefully bounded blur as candidate operations. Exact probabilities and
implementations are controlled within each model experiment. Exact policies
and ablation results belong in the corresponding experiment sections.

\section{Experiments}
\label{sec:experiments}

\subsection{CNN-based model}
\label{subsec:cnn-model}

\subsubsection{Model Overview}
\label{subsubsec:cnn-overview}

The CNN detector is Faster R-CNN \cite{ren2015faster} with a ResNet-50 backbone
\cite{he2016resnet} and an FPN \cite{lin2017fpn}. Its two-stage design first
generates class-agnostic candidates and then performs proposal-level
classification and localization. This is suitable for small and overlapping
objects because proposal quality can be improved separately from the final
detection head.

ResNet-50 extracts hierarchical convolutional features. FPN combines
semantically strong low-resolution maps with spatially detailed high-resolution
maps. The RPN predicts objectness and candidate boxes at every pyramid level;
RoI Align then creates fixed-size proposal features for the classification and
box-regression heads. The classifier predicts background or one of ten
VisDrone classes.

\begin{figure}[htbp]
  \centering
  \includegraphics[width=\textwidth]{\reportfigdir/faster_rcnn_architecture.pdf}
  \caption{Faster R-CNN ResNet-50-FPN architecture used by the CNN experiments.
  The RPN and detection head share the multi-scale features produced by the
  backbone and FPN.}
  \label{fig:faster-rcnn-architecture}
\end{figure}

All runs use Torchvision implementations and COCO pretrained weights
\cite{torchvision_fasterrcnn}. The COCO predictor is replaced by a
\texttt{FastRCNNPredictor} with 11 outputs. F0--F3 use
\texttt{fasterrcnn\_resnet50\_fpn}; F4--F5 use
\texttt{fasterrcnn\_resnet50\_fpn\_v2}. Each experiment starts independently
from COCO weights. The total training loss is

\begin{equation}
L=L_{\mathrm{cls}}+L_{\mathrm{box}}+
L_{\mathrm{rpn\_objectness}}+L_{\mathrm{rpn\_box}}.
\end{equation}

\subsubsection{Set up}
\label{subsubsec:cnn-setup}

Training was performed in Kaggle on an NVIDIA Tesla T4 using CUDA automatic
mixed precision. In the CNN pipeline, the loader converts COCO boxes from
\texttt{xywh} to \texttt{xyxy} and RGB pixels to \texttt{float32} tensors in
$[0,1]$. The Faster R-CNN transform preserves aspect ratio and applies ImageNet
normalization with $\boldsymbol{\mu}=(0.485,0.456,0.406)$ and
$\boldsymbol{\sigma}=(0.229,0.224,0.225)$. F0--F4 resize the shorter side to
800 pixels and cap the longer side at 1,333 pixels; F5 uses 896 and 1,493.
All full runs used 25 epochs, SGD with learning rate 0.0025,
momentum 0.9, weight decay 0.0005, seed 42, two data-loader workers, and
validation batch size one. MultiStepLR reduced the learning rate by 0.1 at
epochs 15 and 20: 0.0025 for epochs 1--14, 0.00025 for epochs 15--19, and
0.000025 for epochs 20--25. Physical batch size and gradient accumulation were
adjusted for memory, but the effective batch remained two images.

\begin{table}[htbp]
  \centering
  \caption{Controlled CNN ablation sequence.}
  \label{tab:cnn-ablation}
  \resizebox{\textwidth}{!}{%
  \begin{tabular}{lllllll}
    \toprule
    Run & Detector & Main intervention & Aug. & Resize & Batch/accum. & Best metric \\
    \midrule
    F0 & V1 & baseline, default anchors/proposals & No  & 800/1,333 & 2/1 & COCO mAP \\
    F1 & V1 & F0 + bbox-safe augmentation        & Yes & 800/1,333 & 2/1 & COCO mAP \\
    F2 & V1 & anchors 8, 16, 32, 64, 128 px     & No  & 800/1,333 & 2/1 & COCO mAP \\
    F3 & V1 & increased crowded-scene proposals  & No  & 800/1,333 & 2/1 & COCO mAP \\
    F4 & V2 & F3 + improved V2 detector           & No  & 800/1,333 & 1/2 & COCO mAP \\
    F5 & V2 & F4 + increased input resolution     & No  & 896/1,493 & 1/2 & VisDrone AP \\
    \bottomrule
  \end{tabular}}
\end{table}

F1 evaluated a CNN-specific online, bounding-box-safe augmentation policy.

\begin{table}[htbp]
  \centering
  \caption{Online augmentation policy evaluated only in CNN experiment F1.}
  \label{tab:augmentation}
  \begin{tabular}{lcl}
    \toprule
    Transformation & Probability & Parameters \\
    \midrule
    Horizontal flip       & 0.50 & left--right \\
    Vertical flip         & 0.20 & top--bottom \\
    Affine transform      & 0.60 & rotation $\pm10^{\circ}$, translation $\pm5\%$, scale 0.90--1.10 \\
    Brightness/contrast   & 0.80 & factors in 0.85--1.15 \\
    Gaussian blur         & 0.10 & radius 0.10--1.00 \\
    \bottomrule
  \end{tabular}
\end{table}

The training script attaches \texttt{ControlledDetectionAugmenter} only for
\texttt{EXPERIMENT=F1}; F0 and F2--F5 use \texttt{augmentation=None}. Images
and box corners are transformed together, affine padding uses RGB
$(114,114,114)$, and boxes are clipped afterward. A box and its synchronized
class label are removed when a coordinate is non-finite, width or height is
below 2 pixels, or retained area is below 30\% of expected post-scale area.
Seed 42 and deterministic worker seeding make the stochastic pipeline
reproducible. F1 otherwise keeps the split, initialization, optimizer,
schedule, effective batch size, resizing, anchors, and proposal limits
identical to F0, and starts independently from the same COCO weights.

Under the unified evaluator, F1 raises AP50 by 1.28 points and AP-small by 0.18
but reduces overall AP by 0.24, AP75 by 0.51, AR@100 by 0.74, and AR@500 by
0.74. The policy was not retained in F2--F5. This is a CNN ablation result, not
a general conclusion that augmentation is harmful to other model families.

F2 used anchor sizes 8, 16, 32, 64, and 128 pixels with aspect ratios 0.5,
1.0, and 2.0 on every FPN level. Relative to F0, it improved AP by 0.54 points,
AP50 by 2.91, AP-small by 1.92, AR-small by 4.88, and AR@500 by 2.65, while
AP75 fell by 0.48. The default anchors were restored for the subsequent
proposal ablation because small anchors improved retrieval more than high-IoU
localization. F3 instead increased RPN pre/post-NMS top-N limits to
4,000/2,000 during training and 2,000/1,000 during inference. Relative to F0,
this improved AP by 1.25, AP75 by 0.98, AR@100 by 2.55, and AR@500 by 2.62
points, so the proposal change was retained.

F4 adopted Faster R-CNN V2. Relative to F3 under the common COCO-style
evaluator, it improved AP by 0.67 points, AP75 by 1.48, AR@100 by 1.26,
AR@500 by 1.22, AP-small by 1.53, and AR-small by 2.22 points, despite a 0.55
point reduction in AP50. F5 then increased input size to 896/1,493.
The attempted 1,024/1,707 setting exceeded Tesla T4 memory. F5 also evaluated
up to 500 detections per image at score threshold 0.001 and selected
\texttt{best.pth} directly by VisDrone AP.

\paragraph{CNN results.}
The saved best checkpoints from all six full runs were re-evaluated on the
same 548 processed validation images with \texttt{pycocotools}. The protocol
uses bounding-box IoU thresholds 0.50:0.05:0.95, 101 recall thresholds, and
\texttt{maxDets=[1,10,100,500]}. AP, AP50, AP75, and every scale-specific
AP/AR value use 500 detections per image; AR@K uses the stated K. COCO area
ranges define small objects below $32^2$ pixels, medium objects from $32^2$ to
$96^2$, and large objects at least $96^2$ pixels. Predictions are retained at
score 0.001 and capped at 500 per image.

This is a COCO-style extension to maxDet=500, not the official COCO maxDet=100
summary or the official VisDrone ignored-region evaluator. The processed
ground truth excludes ignored rows, so these results support controlled project
comparisons and must not be mixed directly with official-compatible VisDrone
scores.

\begin{table}[htbp]
  \centering
  \caption{Unified COCO-style validation results at maxDet=500: overall precision and recall (\%).}
  \label{tab:cnn-results}
  \resizebox{\textwidth}{!}{%
  \begin{tabular}{lrrrrrrr}
    \toprule
    Run & AP (50:95) & AP50 & AP75 & AR@1 & AR@10 & AR@100 & AR@500 \\
    \midrule
    F0 & 23.972 & 42.603 & 23.507 & 10.490 & 30.821 & 40.604 & 40.618 \\
    F1 & 23.733 & 43.881 & 22.998 & 10.403 & 30.306 & 39.861 & 39.879 \\
    F2 & 24.508 & 45.513 & 23.028 & 10.424 & 30.931 & 43.029 & 43.272 \\
    F3 & 25.223 & 45.201 & 24.482 & 10.787 & 31.618 & 43.150 & 43.237 \\
    F4 & 25.893 & 44.652 & 25.958 & 10.878 & 32.831 & 44.407 & 44.459 \\
    \textbf{F5} & \textbf{27.121} & \textbf{47.308} & \textbf{26.806} & \textbf{11.300} & \textbf{34.411} & \textbf{47.024} & \textbf{47.114} \\
    \bottomrule
  \end{tabular}}
\end{table}

\begin{table}[htbp]
  \centering
  \caption{Unified COCO-style validation results at maxDet=500 by object area (\%).}
  \label{tab:cnn-scale-results}
  \resizebox{\textwidth}{!}{%
  \begin{tabular}{lrrrrrr}
    \toprule
    Run & AP small & AR small & AP medium & AR medium & AP large & AR large \\
    \midrule
    F0 & 15.280 & 31.655 & 34.305 & 52.893 & 42.682 & 58.579 \\
    F1 & 15.455 & 30.996 & 33.504 & 52.907 & 44.308 & 61.690 \\
    F2 & 17.199 & 36.539 & 33.156 & 51.572 & 42.101 & 57.282 \\
    F3 & 16.871 & 35.238 & 35.192 & 53.994 & 43.417 & 57.516 \\
    F4 & 18.397 & 37.453 & 35.577 & 53.527 & 43.221 & 55.767 \\
    \textbf{F5} & \textbf{18.954} & \textbf{40.025} & \textbf{37.783} & \textbf{56.683} & \textbf{48.531} & \textbf{65.439} \\
    \bottomrule
  \end{tabular}}
\end{table}

\begin{figure}[htbp]
  \centering
  \includegraphics[width=\textwidth]{\reportfigdir/cnn_experiment_progression.pdf}
  \caption{Unified COCO-style precision and recall at maxDet=500. F1 shows
  that augmentation alone did not improve overall AP, while F2--F5 improve
  the baseline and F5 achieves the best value for every reported metric.}
  \label{fig:cnn-progression}
\end{figure}

F5 is the final CNN model. Compared with F4, it gains 1.228 points AP, 2.656
AP50, 0.848 AP75, 2.617 AR@100, and 2.655 AR@500. It also gains 0.557
AP-small, 2.206 AP-medium, and 5.310 AP-large. This AP gain is 4.74\% in
relative terms. Compared with F0, F5 gains 3.149 AP points (13.14\% relative),
4.705 AP50, 3.299 AP75, 6.496 AR@500, and 3.674 AP-small. Training required
9.36 hours and peaked at 11.28 GB GPU memory. Its best checkpoint under the
training-time VisDrone criterion occurred at epoch 5: although training loss
fell from 1.1165 to 0.4009, that validation AP declined to 26.15\% by epoch 25.
Final inference must therefore use \texttt{f5/best.pth}, not
\texttt{last.pth}.

\begin{figure}[htbp]
  \centering
  \includegraphics[width=0.92\textwidth]{\reportfigdir/f5_training_dynamics.pdf}
  \caption{F5 training dynamics. Training loss decreases throughout the run,
  whereas validation AP reaches its maximum at epoch 5 and subsequently
  fluctuates. This supports validation-based checkpoint selection.}
  \label{fig:f5-training-dynamics}
\end{figure}

A major remaining limitation is small-object detection: F5 AP-small is
18.954\%, versus 37.783\% AP-medium and 48.531\% AP-large. Under the separate
official-compatible VisDrone protocol, a 640$\times$640 tiled-inference
experiment increased AR@500 by 0.77 points but reduced AP by 0.88 and AP75 by
1.23 points. It was rejected as the default method. Other limitations include
crowded scenes and class imbalance; future work should test conservative small-object-aware
augmentation, class-aware optimization, and multiple random seeds.
